In this lecture, we state the law of large numbers and return to random walks on graphs. We describe
a random walk by its transition matrix $P$, show that the distribution of the token evolves as
$p_{t+1} = p_t P$, find the stationary distribution $\pi = \pi P$, and prove by an analysis of the
eigenvalues of $P$ that the walk on a connected, non-bipartite graph converges to it.

\section{Law of Large Numbers}

Let $X_1, X_2, \dots, X_n$ be independent, identically distributed (i.i.d.) random variables, each
distributed as a random variable $X$ with finite expectation $\expectation{X}$. The law of large numbers says that the average of the samples approaches
the expectation.

\begin{theorem}[Law of Large Numbers, as Stated in Class]\label{thm:slln-0910}
    Let $X_1, X_2, \dots$ be i.i.d.\ random variables distributed as $X$, where $\expectation{X}$ is
    finite. Then, with probability $1$,
    \[\lim_{n \to \infty} \frac{\sum_{i=1}^{n} X_i}{n} = \expectation{X}.\]
\end{theorem}

\paragraph{Remark.} There are two laws of large numbers, and they differ in the kind of convergence.
The statement above, convergence with probability $1$, is the \emph{strong} law; we do not prove it
here. The \emph{weak} law is convergence in probability, which is implied by convergence with
probability $1$. The weak law is the one that we proved from Chebyshev's inequality in
\cref{thm:wlln-0908}, under the additional assumption that $\variance{X}$ is finite. In the notes the
heading is ``weak law'' and the statement written under it is the strong one.

\section{Random Walks on Graphs}

As in the last lecture, a \emph{random walk} on a finite undirected graph $G = (V, E)$ moves a token in
every step from its current vertex $u$ to a uniformly random neighbour of $u$. We take
$V = \set{1, \dots, n}$, write $m = \abs{E}$ for the number of edges, $u \sim v$ if $\set{u,v} \in E$, and
$\deg(u)$ (or $d_u$) for the number of neighbours of $u$, and we assume throughout that $G$ has no
isolated vertices. Let $X_t$ be the vertex of the token at time $t = 0, 1, 2, \dots$.

\begin{definition}[Distribution at Time $t$, Transition Matrix]\label{def:walk-0910}
    The \textbf{distribution at time $t$} is the row vector
    \[p_t = \left(p_t^1, p_t^2, \dots, p_t^n\right), \qquad p_t^i = \prob{X_t = i},\]
    i.e., $p_t^i$ is the probability that the token is at vertex $i$ at time $t$ (the subscript is the
    time, the superscript is the vertex). The \textbf{transition matrix} is the $n \times n$ matrix $P$
    with entries
    \[P_{uv} = \prob{u \to v} = \prob{X_{t+1} = v \given X_t = u}
        = \begin{cases}
            \frac{1}{\deg(u)} & \text{if } u \sim v, \\
            0                 & \text{otherwise.}
        \end{cases}\]
\end{definition}

The entries $P_{uv}$ do not depend on $t$, and the next vertex depends only on the current one. A process
with these two properties is a (time-homogeneous) \emph{Markov chain}.

\subsection{The Triangle}

Consider the triangle on the vertices $1, 2, 3$ with the token starting at vertex $1$
(the right half of \cref{fig:walks-0908}). Every vertex has degree $2$, so the token moves to each of the two other
vertices with probability $\sfrac{1}{2}$ and
\[P = \begin{pmatrix}
        0            & \sfrac{1}{2} & \sfrac{1}{2} \\
        \sfrac{1}{2} & 0            & \sfrac{1}{2} \\
        \sfrac{1}{2} & \sfrac{1}{2} & 0
    \end{pmatrix}.\]

\subsection{The Evolution \texorpdfstring{$p_{t+1} = p_t P$}{p(t+1) = p(t) P}}

\begin{lemma}[Evolution of the Distribution]\label{lem:evolution-0910}
    For every $t \geq 0$,
    \[p_{t+1} = p_t P, \qquad \text{and hence} \qquad p_t = p_0 P^t.\]
\end{lemma}

\begin{proof}
    Fix a vertex $v$. The events $\set{X_t = u}$, $u \in V$, are disjoint and one of them occurs, so
    \[p_{t+1}^v = \prob{X_{t+1} = v} = \sum_{u \in V} \prob{X_t = u \text{ and } X_{t+1} = v}
        = \sum_{u \in V} p_t^u P_{uv} = (p_t P)_v.\]
    Here $\prob{X_t = u \text{ and } X_{t+1} = v} = p_t^u P_{uv}$ by the definition of $P_{uv}$ as a
    conditional probability when $p_t^u > 0$, and both sides are $0$ when $p_t^u = 0$. The last expression
    is the $\nth{v}$ entry of the product of the row vector $p_t$ with the matrix $P$.
    This holds for every $v$, so $p_{t+1} = p_t P$. The formula $p_t = p_0 P^t$ follows by induction on $t$:
    it is true for $t = 0$ since $P^0 = I$, and if $p_t = p_0 P^t$ then $p_{t+1} = p_t P = p_0 P^{t+1}$.
\end{proof}

For the triangle we can verify this directly. Since $p_0 = (1,0,0)$, the product $p_0 P$ is the first row
of $P$, and
\begin{align*}
    p_1 & = p_0 P = \left(0, \sfrac{1}{2}, \sfrac{1}{2}\right),                                                 \\
    p_2 & = p_1 P = p_0 P^2
    = \tfrac{1}{2}\left(\sfrac{1}{2}, 0, \sfrac{1}{2}\right) + \tfrac{1}{2}\left(\sfrac{1}{2}, \sfrac{1}{2}, 0\right)
    = \left(\sfrac{1}{2}, \sfrac{1}{4}, \sfrac{1}{4}\right),
\end{align*}
which are the vectors computed in the last lecture. The closed form of $p_t$ for all $t$ is \cref{lem:triangle-0908}.

\subsection{Row Stochastic Matrices}

Every row of the triangle's matrix $P$ has nonnegative entries that sum to $1$; in other words, the $L_1$
norm of every row is $1$.

\begin{definition}[Probability Vector, Row Stochastic Matrix]\label{def:stochastic-0910}
    A vector $p = (p^1, \dots, p^n)$ is a \textbf{probability vector} if $p^u \geq 0$ for all $u$ and
    $\sum_{u=1}^{n} p^u = 1$. An $n \times n$ matrix $P$ is \textbf{row stochastic} if every row of $P$ is a
    probability vector:
    \[P_{uv} \geq 0 \text{ for all } u, v, \qquad \text{and} \qquad \sum_{v=1}^{n} P_{uv} = 1 \text{ for all } u.\]
\end{definition}

The transition matrix of a random walk on any graph is row stochastic: its entries are $0$ or
$\sfrac{1}{\deg(u)}$, and row $u$ has exactly $\deg(u)$ entries equal to $\sfrac{1}{\deg(u)}$. The
nonnegativity is part of the definition and is needed below.

% Page 2, bottom, standalone line: "Show that p_{t+1} is a probability vector, assume p_t is."
\needspace{9\baselineskip} % keep the exercise and the start of its solution on one page
\begin{exercise}\label{ex:probvec-0910}
    Show that $p_{t+1}$ is a probability vector, assuming that $p_t$ is.
\end{exercise}

The solution is \cref{thm:probvec-0910}.

\begin{theorem}[A Row Stochastic Matrix Preserves Probability Vectors]\label{thm:probvec-0910}
    Let $P$ be a row stochastic $n \times n$ matrix and let $p$ be a probability vector. Then $pP$ is a
    probability vector. Consequently, if $p_0$ is a probability vector and $p_{t+1} = p_t P$ for all
    $t \geq 0$, then $p_t = p_0 P^t$ is a probability vector for every $t \geq 0$.
\end{theorem}

\begin{proof}
    Let $q = pP$, so that $q^v = \sum_{u=1}^{n} p^u P_{uv}$ for every $v$.

    \emph{Nonnegativity.} Every term $p^u P_{uv}$ is a product of two nonnegative numbers, since $p^u \geq 0$
    ($p$ is a probability vector) and $P_{uv} \geq 0$ ($P$ is row stochastic). A sum of nonnegative numbers is
    nonnegative, so $q^v \geq 0$ for every $v$.

    \emph{The entries sum to $1$.} We have
    \[\sum_{v=1}^{n} q^v = \sum_{v=1}^{n} \sum_{u=1}^{n} p^u P_{uv}
        = \sum_{u=1}^{n} \sum_{v=1}^{n} p^u P_{uv}
        = \sum_{u=1}^{n} p^u \left(\sum_{v=1}^{n} P_{uv}\right)
        = \sum_{u=1}^{n} p^u \cdot 1 = 1.\]
    The second equality exchanges the order of two finite sums (both sides add up the same $n^2$
    numbers $p^u P_{uv}$), the fourth uses that row $u$ of $P$ sums to $1$, and the last that the entries
    of $p$ sum to $1$. Hence $q = pP$ is a probability vector.

    \emph{All times.} We prove by induction on $t$ that $p_t$ is a probability vector. For $t = 0$ this is
    the assumption on $p_0$. If $p_t$ is a probability vector, then by the first part (with $p = p_t$) so is
    $p_t P = p_{t+1}$, which is exactly the statement of \cref{ex:probvec-0910}.
\end{proof}

Both halves of ``row stochastic'' are used: the nonnegativity of the entries of $P$ gives $p_{t+1}^v \geq 0$,
and the unit row sums give $\sum_v p_{t+1}^v = 1$.

\subsection{The Transition Matrix of a General Graph}

Let $A$ be the \emph{adjacency matrix} of $G$, i.e., $A_{uv} = 1$ if $u \sim v$ and $A_{uv} = 0$ otherwise,
and let $D$ be the diagonal matrix of the degrees. Then
\[D = \begin{pmatrix}
        d_1 &        &     \\
            & \ddots &     \\
            &        & d_n
    \end{pmatrix},
    \qquad
    D^{-1} = \begin{pmatrix}
        \sfrac{1}{d_1} &        &                \\
                       & \ddots &                \\
                       &        & \sfrac{1}{d_n}
    \end{pmatrix},
    \qquad
    P = D^{-1} A.\]
Indeed, multiplying $A$ from the left by the diagonal matrix $D^{-1}$ multiplies row $u$ of $A$ by
$\sfrac{1}{d_u}$, so $(D^{-1} A)_{uv} = \sfrac{A_{uv}}{\deg(u)}$, which is $P_{uv}$ of
\cref{def:walk-0910}. The inverse $D^{-1}$ exists because there are no isolated vertices. (The diagonal
matrix drawn in the notes next to $D^{-1}$ shows the entries $d_1, \dots, d_n$, i.e., $D$ itself.)

If $G$ is \emph{$d$-regular}, i.e., every vertex has the same degree $d$, then $D = d I$ and the formula
becomes
\[P = \frac{A}{d}.\]
This is the regular case only; the triangle is $2$-regular, and its $P$ above is $\sfrac{A}{2}$.

\section{Stationary Distribution}

\begin{definition}[Stationary Distribution]\label{def:stationary-0910}
    A probability vector $\pi$ is a \textbf{stationary distribution} of the walk with transition matrix $P$
    if
    \[\pi P = \pi \cdot 1 = \pi,\]
    i.e., $\pi$ is a left eigenvector of $P$ with eigenvalue $1$.
\end{definition}

\begin{figure}[ht]
    \centering
    \begin{tikzpicture}
        \node[point] (u) at (0,0) {};
        \node[point] (v) at (2,1.3) {};
        \node[point] (w1) at (-1.4,-1.2) {};
        \node[point] (w2) at (1.4,-1.2) {};
        \draw (u) -- (v);
        \draw (u) -- (w1);
        \draw (u) -- (w2);
        \node[left=3pt] at (u) {$u$};
        \node[above=3pt] at (v) {$v$};
        \node[right] at (1.3,0.45) {\small $P_{vu} = \sfrac{1}{\deg(v)}$};
    \end{tikzpicture}
    \caption{A vertex $u$ and its neighbours. Each neighbour $v$ contributes $\sfrac{\pi_v}{\deg(v)}$ to
        $(\pi P)_u$.}
    \label{fig:star-0910}
\end{figure}

By \cref{lem:evolution-0910}, if $p_t = \pi$ then $p_{t+1} = \pi P = \pi$: a walk whose distribution is
$\pi$ keeps the distribution $\pi$ forever. To find $\pi$ we write out the $\nth{u}$ entry of $\pi P$.
Only the neighbours $v$ of $u$ have $P_{vu} \neq 0$, and for them $P_{vu} = \sfrac{1}{\deg(v)}$
(\cref{fig:star-0910}), so
\begin{equation}\label{eq:piP-0910}
    (\pi P)_u = \sum_{v \in V} \pi_v P_{vu} = \sum_{v \sim u} \frac{\pi_v}{\deg(v)}.
\end{equation}
(The notes stop at $(\pi P)_u = \sum_{v \sim u} \pi_v$, in which the factor $\sfrac{1}{\deg(v)}$ is
missing.) In words, the probability of being at $u$ after one step is collected from the neighbours of
$u$, and a neighbour $v$ sends to $u$ a $\sfrac{1}{\deg(v)}$ fraction of its own probability.

\begin{theorem}[Stationary Distribution of a Random Walk on a Graph]\label{thm:stationary-0910}
    Let $G$ be a graph with $m$ edges and no isolated vertices. Then
    \[\pi_v = \frac{\deg(v)}{2m} \qquad (v \in V)\]
    is a stationary distribution of the random walk on $G$.
\end{theorem}

\begin{proof}
    \emph{$\pi$ is a probability vector.} Since there are no isolated vertices and $V$ is not empty, we
    have $m \geq 1$, so $\pi$ is well defined, and $\pi_v \geq 0$. Every edge $\set{u,v}$ is counted exactly
    twice in $\sum_{v \in V} \deg(v)$, once in $\deg(u)$ and once in $\deg(v)$, so
    $\sum_{v \in V} \deg(v) = 2m$ (the handshake lemma). Hence $\sum_{v \in V} \pi_v = \sfrac{2m}{2m} = 1$.

    \emph{$\pi P = \pi$.} For every vertex $u$, by \cref{eq:piP-0910},
    \[(\pi P)_u = \sum_{v \sim u} \frac{\pi_v}{\deg(v)}
        = \sum_{v \sim u} \frac{\deg(v)}{2m} \cdot \frac{1}{\deg(v)}
        = \sum_{v \sim u} \frac{1}{2m}
        = \frac{\deg(u)}{2m} = \pi_u,\]
    where we used $\deg(v) \geq 1$ to cancel, and that the sum has one term for each of the $\deg(u)$
    neighbours of $u$.
\end{proof}

For a $d$-regular graph on $n$ vertices we have $2m = nd$, so $\pi_v = \sfrac{d}{nd} = \sfrac{1}{n}$: the
stationary distribution is uniform. For the triangle ($m = 3$, all degrees $2$) it is
$\pi = \left(\sfrac{1}{3}, \sfrac{1}{3}, \sfrac{1}{3}\right)$, which is the limit of the closed form of
\cref{lem:triangle-0908} as $t \to \infty$.

\section{Convergence to the Stationary Distribution}

The reason to study $\pi$ is that, as $t \to \infty$, the distribution $p_t$ of the walk approaches it:
$p_t \to \pi$. Here convergence of vectors is entry by entry. Two conditions on the graph were given in
class: (1) the graph should be connected, and (2) it should be non-bipartite. These are the conditions
for the walk to converge to one and the same limit from every start; a stationary distribution itself
exists for every graph without isolated vertices, by \cref{thm:stationary-0910}.

% Pages 3-4: "As t -> oo, p_t -> pi" with the two conditions (connected, non-bipartite). Assigned in
% class: prove the convergence by eigenvalue analysis.
\needspace{9\baselineskip} % keep the exercise and the start of its solution on one page
\begin{exercise}\label{ex:convergence-0910}
    Let $G$ be connected and non-bipartite. Prove, by analysing the eigenvalues of the transition
    matrix $P$, that the random walk on $G$ converges to the stationary distribution: $p_t \to \pi$ as
    $t \to \infty$, for every starting distribution $p_0$.
\end{exercise}

The solution is \cref{thm:convergence-0910}, which is proved in \cref{sec:eigen-0910}.

\begin{theorem}[Convergence of Random Walks]\label{thm:convergence-0910}
    Let $G$ be a connected, non-bipartite graph with $m$ edges. Then for every starting distribution
    $p_0$,
    \[p_t = p_0 P^t \longrightarrow \pi \qquad \text{as } t \to \infty, \qquad \text{where } \pi_v = \frac{\deg(v)}{2m}.\]
    In particular, $\pi$ is the only stationary distribution of the walk.
\end{theorem}

\subsection{Proof by Eigenvalue Analysis}\label{sec:eigen-0910}

The idea is to write $P^t$ in terms of the eigenvalues of $P$. The eigenvalue $1$ gives the limit $\pi$,
and every other eigenvalue $\lambda$ has $\abs{\lambda} < 1$, so its contribution $\lambda^t$ dies out.
Connectedness is what makes the eigenvalue $1$ occur only once, and non-bipartiteness is what excludes
the eigenvalue $-1$, whose powers $(-1)^t$ would not die out.

In this subsection $p_t$ and $\pi$ are row vectors as before, while $x$, $w$ and $z$ denote real
column vectors with entries $x_u$ ($u \in V$), and $\mathbf{1}$ is the column vector with all entries $1$.
A nonzero column vector $x$ is an \emph{eigenvector} of a matrix $M$ with \emph{eigenvalue} $\lambda$ if
$Mx = \lambda x$.

\paragraph{A symmetric matrix with the same eigenvalues.} The adjacency matrix $A$ is symmetric, but
$P = D^{-1} A$ is symmetric only if all degrees are equal. Let $D^{1/2}$ be the diagonal matrix with
entries $\sqrt{d_1}, \dots, \sqrt{d_n}$ and $D^{-1/2}$ its inverse, and put
\[N = D^{-1/2} A D^{-1/2}, \qquad \text{i.e.,} \qquad N_{uv} = \frac{A_{uv}}{\sqrt{d_u d_v}}.\]
Since $A_{uv} = A_{vu}$, the matrix $N$ is symmetric. Moreover,
\[D^{-1/2} N D^{1/2} = D^{-1/2} D^{-1/2} A D^{-1/2} D^{1/2} = D^{-1} A = P,\]
and in the product of $t$
copies of $D^{-1/2} N D^{1/2}$ every inner pair $D^{1/2} D^{-1/2}$ cancels, so
\begin{equation}\label{eq:similar-0910}
    P = D^{-1/2} N D^{1/2}, \qquad P^t = D^{-1/2} N^t D^{1/2} \quad \text{for all } t \geq 0.
\end{equation}
Eigenvectors of $N$ give eigenvectors of $P$ with the same eigenvalue: if $Nw = \lambda w$ and
$w \neq 0$, then $x = D^{-1/2} w \neq 0$ and
\begin{equation}\label{eq:transfer-0910}
    P x = D^{-1/2} N D^{1/2} D^{-1/2} w = D^{-1/2} N w = \lambda D^{-1/2} w = \lambda x.
\end{equation}

\paragraph{The spectral theorem.} For symmetric matrices we use the following standard theorem of
linear algebra, without proof.

\begin{theorem}[Spectral Theorem for Real Symmetric Matrices]\label{thm:spectral-0910}
    Let $N$ be a real symmetric $n \times n$ matrix. Then there are real numbers
    $\lambda_1, \dots, \lambda_n$ and real column vectors $w_1, \dots, w_n$ such that
    $N w_i = \lambda_i w_i$ for all $i$, and $w_1, \dots, w_n$ is an orthonormal basis of
    $\mathbb{R}^n$: $w_i^{T} w_i = 1$, and $w_i^{T} w_j = 0$ for $i \neq j$.
\end{theorem}

Let $x$ be a real column vector. Since $w_1, \dots, w_n$ is a basis, $x = \sum_{i} c_i w_i$ for unique
numbers $c_i$, and multiplying from the left by $w_j^{T}$ gives $c_j = w_j^{T} x$. Applying $N$ to this
expansion $t$ times multiplies the $\nth{i}$ term by $\lambda_i^t$, so
\[N^t x = \sum_{i=1}^{n} \lambda_i^t \left(w_i^{T} x\right) w_i
    = \left(\sum_{i=1}^{n} \lambda_i^t\, w_i w_i^{T}\right) x.\]
This holds for every $x$, so
\begin{equation}\label{eq:powers-0910}
    N^t = \sum_{i=1}^{n} \lambda_i^t\, w_i w_i^{T} \qquad \text{for all } t \geq 0,
\end{equation}
where $w_i w_i^{T}$ is the $n \times n$ matrix with entries $(w_i)_u (w_i)_v$.

\paragraph{The eigenvalues of $P$.} For a column vector $x$, the $\nth{u}$ entry of $Px$ is the average
of $x$ over the neighbours of $u$:
\begin{equation}\label{eq:average-0910}
    (Px)_u = \sum_{v \in V} P_{uv} x_v = \frac{1}{\deg(u)} \sum_{v \sim u} x_v.
\end{equation}

\begin{lemma}[Eigenvalues of the Transition Matrix]\label{lem:eigen-0910}
    Let $G$ be a connected graph with at least two vertices, let $\lambda$ be a real number, and let
    $x \neq 0$ be a real column vector with $Px = \lambda x$.
    \begin{enumerate}
        \item $\abs{\lambda} \leq 1$.
        \item If $\lambda = 1$, then all entries of $x$ are equal, i.e., $x = c \mathbf{1}$ for some $c \neq 0$.
        \item If $\lambda = -1$, then $G$ is bipartite.
    \end{enumerate}
\end{lemma}

\begin{proof}
    A connected graph with at least two vertices has no isolated vertices, so $P$ is defined and every
    vertex has a neighbour. Let $M = \max_{u} \abs{x_u}$; since $x \neq 0$, we have $M > 0$.

    \emph{Part 1.} Let $u$ be a vertex with $\abs{x_u} = M$. By \cref{eq:average-0910} and the triangle
    inequality,
    \[\abs{\lambda} M = \abs{\lambda x_u} = \abs{(Px)_u}
        \leq \frac{1}{\deg(u)} \sum_{v \sim u} \abs{x_v} \leq \frac{1}{\deg(u)} \cdot \deg(u) \cdot M = M.\]
    Dividing by $M > 0$ gives $\abs{\lambda} \leq 1$.

    \emph{Part 2.} Let $\lambda = 1$, let $K = \max_u x_u$, and let $S = \set{u \suchthat x_u = K}$, which
    is not empty. Let $u \in S$. By \cref{eq:average-0910}, $K = x_u = (Px)_u$ is the average of the
    numbers $x_v$ over the neighbours $v$ of $u$. Each of these numbers is at most $K$, and if one of
    them were smaller than $K$ the average would be smaller than $K$. So $x_v = K$ for every neighbour
    $v$ of $u$: every neighbour of a vertex of $S$ is in $S$. Since $G$ is connected, every vertex can
    be reached from a vertex of $S$ along a path, so $S = V$ and $x = K \mathbf{1}$. Finally $K \neq 0$
    because $x \neq 0$.

    \emph{Part 3.} Let $\lambda = -1$, and let
    \[S_+ = \set{u \suchthat x_u = M}, \qquad S_- = \set{u \suchthat x_u = -M}.\]
    These two sets are disjoint because $M > 0$, and at least one of them is not empty. Let $u \in S_+$.
    By \cref{eq:average-0910}, the average of $x_v$ over the neighbours $v$ of $u$ is
    $(Px)_u = -x_u = -M$. Each $x_v$ is at least $-M$, so, as in Part 2, all of them are equal to
    $-M$: every neighbour of a vertex of $S_+$ is in $S_-$. In the same way, if $u \in S_-$ then the
    average of $x_v$ over the neighbours of $u$ is $-x_u = M$, each $x_v$ is at most $M$, and so every
    neighbour of a vertex of $S_-$ is in $S_+$. Hence every neighbour of a vertex of $S_+ \cup S_-$ is
    again in $S_+ \cup S_-$, and since $G$ is connected, $S_+ \cup S_- = V$. So every vertex is in
    exactly one of $S_+$ and $S_-$, and every edge joins a vertex of $S_+$ to a vertex of $S_-$. This
    is a bipartition of $G$.
\end{proof}

\begin{proof}[Proof of \cref{thm:convergence-0910}]
    A non-bipartite graph contains an odd cycle, so $G$ has at least three vertices, and it is
    connected, so \cref{lem:eigen-0910} applies. Let $\lambda_1, \dots, \lambda_n$ and $w_1, \dots, w_n$
    be the eigenvalues and the orthonormal eigenvectors of the symmetric matrix $N$ given by
    \cref{thm:spectral-0910}. By \cref{eq:transfer-0910}, $x_i = D^{-1/2} w_i$ is a real eigenvector of $P$
    with eigenvalue $\lambda_i$, so \cref{lem:eigen-0910} applies to every pair $\lambda_i$, $x_i$.

    \emph{Step 1: the eigenvalue $1$ occurs exactly once.} Let $z = D^{1/2} \mathbf{1}$, the column
    vector with entries $\sqrt{d_u}$. The $\nth{u}$ entry of $A \mathbf{1}$ is the number of neighbours of
    $u$, so $A \mathbf{1} = D \mathbf{1}$ and
    \[N z = D^{-1/2} A D^{-1/2} D^{1/2} \mathbf{1} = D^{-1/2} A \mathbf{1} = D^{-1/2} D \mathbf{1} = D^{1/2} \mathbf{1} = z.\]
    So $z$ is an eigenvector of $N$ with eigenvalue $1$, and
    $z^{T} z = \sum_{u} d_u = 2m$ by the handshake lemma.

    At least one $\lambda_i$ equals $1$: write $z = \sum_i c_i w_i$. Then
    $\sum_i c_i w_i = z = Nz = \sum_i c_i \lambda_i w_i$, and the coefficients with respect to a basis are
    unique, so $c_i = c_i \lambda_i$ for all $i$. Since $z \neq 0$, some $c_i$ is not $0$, and for this
    $i$ we get $\lambda_i = 1$.

    If $\lambda_i = 1$, then by part 2 of \cref{lem:eigen-0910} we have $D^{-1/2} w_i = c \mathbf{1}$ for
    some $c \neq 0$, i.e., $w_i = c D^{1/2} \mathbf{1} = c z$. Hence at most one $\lambda_i$ equals $1$: if
    $\lambda_i = \lambda_j = 1$ with $i \neq j$, then $w_i = c z$ and $w_j = c' z$ with $c, c' \neq 0$,
    and $w_i^{T} w_j = c c' \cdot 2m \neq 0$, which contradicts the orthogonality. We number the
    eigenvalues so that $\lambda_1 = 1$. Then $w_1 = c z$ with $1 = w_1^{T} w_1 = c^2 \cdot 2m$, and
    therefore
    \begin{equation}\label{eq:first-0910}
        w_1 w_1^{T} = c^2 z z^{T} = \frac{1}{2m} z z^{T}.
    \end{equation}

    \emph{Step 2: all other eigenvalues are smaller than $1$ in absolute value.} Let $i \geq 2$. Then
    $\abs{\lambda_i} \leq 1$ by part 1 of \cref{lem:eigen-0910}, $\lambda_i \neq 1$ by Step 1, and
    $\lambda_i \neq -1$ by part 3, because $G$ is not bipartite. Hence
    \[\rho := \max_{2 \leq i \leq n} \abs{\lambda_i} < 1.\]

    \emph{Step 3: the limit of $N^t$.} By \cref{eq:powers-0910,eq:first-0910},
    \[N^t = \frac{1}{2m} z z^{T} + E_t, \qquad \text{where } E_t = \sum_{i=2}^{n} \lambda_i^t\, w_i w_i^{T}.\]
    Every entry of a vector of length $1$ is at most $1$ in absolute value, so every entry of
    $w_i w_i^{T}$ is at most $1$ in absolute value, and therefore
    \[\abs{(E_t)_{uv}} \leq \sum_{i=2}^{n} \abs{\lambda_i}^t \leq (n-1) \rho^t \qquad \text{for all } u, v.\]

    \emph{Step 4: the limit of $p_t$.} By \cref{lem:evolution-0910} and \cref{eq:similar-0910},
    \[p_t = p_0 P^t = p_0 D^{-1/2} N^t D^{1/2}
        = \frac{1}{2m}\, p_0 D^{-1/2} z z^{T} D^{1/2} + p_0 D^{-1/2} E_t D^{1/2}.\]
    In the first term, $D^{-1/2} z = \mathbf{1}$, and $p_0 \mathbf{1} = \sum_u p_0^u = 1$ because $p_0$ is a
    probability vector; and $z^{T} D^{1/2} = \mathbf{1}^{T} D = (d_1, \dots, d_n)$. So the first term is
    \[\frac{1}{2m} \left(p_0 \mathbf{1}\right) \left(z^{T} D^{1/2}\right) = \frac{1}{2m} (d_1, \dots, d_n) = \pi,\]
    for every $t$ and every $p_0$. The $\nth{v}$ entry of the second term is
    $\sum_{u} p_0^u\, d_u^{-1/2}\, (E_t)_{uv}\, d_v^{1/2}$. Since $1 \leq d_u$ and $d_v \leq n$, we have
    $\sqrt{d_v / d_u} \leq \sqrt{n}$, and with Step 3 and $\sum_u p_0^u = 1$ we obtain
    \begin{equation}\label{eq:rate-0910}
        \abs{p_t^v - \pi_v} \leq \sum_{u} p_0^u \sqrt{\frac{d_v}{d_u}}\, \abs{(E_t)_{uv}}
        \leq \sqrt{n}\, (n-1)\, \rho^t \leq n^{3/2} \rho^t \qquad \text{for every } v \in V.
    \end{equation}
    Since $0 \leq \rho < 1$, the right-hand side tends to $0$ as $t \to \infty$. Hence $p_t \to \pi$.

    \emph{Uniqueness.} Let $\pi'$ be any stationary distribution, and start the walk with $p_0 = \pi'$.
    Then $p_t = \pi'$ for all $t$, and $p_t \to \pi$ forces $\pi' = \pi$.
\end{proof}

\paragraph{The rate of convergence.} \Cref{eq:rate-0910} says more than the theorem: the distance to
$\pi$ is at most a constant times $\rho^t$, i.e., it decreases geometrically, where $\rho$ is the largest absolute
value of an eigenvalue of $P$ other than the eigenvalue $1$. The closer $\rho$ is to $1$, the slower the
walk converges.

\paragraph{The regular case.} If $G$ is $d$-regular, then $D = dI$ and $N = \sfrac{A}{d} = P$: the
transition matrix itself is symmetric, and no change of matrix is needed. Here $z = \sqrt{d}\, \mathbf{1}$,
and we may take $w_1 = \sfrac{\mathbf{1}}{\sqrt{n}}$. Writing the starting distribution as a column vector
$p_0^{T} = \sum_i c_i w_i$, we have $c_1 = w_1^{T} p_0^{T} = \sfrac{1}{\sqrt{n}}$, and since $P^{T} = P$,
\[p_t^{T} = P^t p_0^{T} = \sum_{i=1}^{n} c_i \lambda_i^t w_i
    = \frac{1}{n} \mathbf{1} + \sum_{i=2}^{n} c_i \lambda_i^t w_i
    \longrightarrow \frac{1}{n} \mathbf{1},\]
the uniform distribution, because $\abs{\lambda_i} < 1$ for $i \geq 2$ by Steps 1 and 2 of the proof.

\paragraph{The triangle again.} The triangle is $2$-regular with $P = \frac{1}{2}(J - I)$, where $J$ is
the $3 \times 3$ matrix with all entries $1$. We have $J \mathbf{1} = 3 \cdot \mathbf{1}$, and $Jx = 0$ for
every $x$ whose entries sum to $0$, for example $x = (1, -1, 0)^{T}$ and $x = (1, 1, -2)^{T}$. So
$P \mathbf{1} = \mathbf{1}$ and $Px = -\frac{1}{2} x$ for these two vectors: the eigenvalues of $P$ are
$1, -\sfrac{1}{2}, -\sfrac{1}{2}$, and $\rho = \sfrac{1}{2}$. For the start $p_0 = (1, 0, 0)$ we have
$p_0^{T} = \frac{1}{3} \mathbf{1} + y$ with $y = \left(\sfrac{2}{3}, -\sfrac{1}{3}, -\sfrac{1}{3}\right)^{T}$,
whose entries sum to $0$, so
\[p_t^{T} = P^t p_0^{T} = \frac{1}{3} \mathbf{1} + \left(-\frac{1}{2}\right)^{t} y,\]
which is exactly the closed form of \cref{lem:triangle-0908}.

\paragraph{Remark.} The theorem is a special case of the convergence theorem for finite Markov chains
that are irreducible and aperiodic; see Levin, Peres and Wilmer, \emph{Markov Chains and Mixing Times},
Chapter~4. A random walk on a graph is irreducible exactly when the graph is connected, and for a
connected graph it is aperiodic exactly when the graph is not bipartite.

\subsection{Both Conditions Are Necessary}

The two graphs of \cref{fig:counter-0910} show that neither condition of \cref{thm:convergence-0910}
can be dropped.

\begin{figure}[ht]
    \centering
    \begin{tikzpicture}
        % two disjoint triangles
        \foreach \sh/\la/\lb/\lc in {0/1/2/3, 3.2/4/5/6} {
                \begin{scope}[xshift=\sh cm]
                    \node[point] (a) at (0,0) {};
                    \node[point] (b) at (1,1.4) {};
                    \node[point] (c) at (2,0) {};
                    \draw (a) -- (b) -- (c) -- (a);
                    \node[below=3pt] at (a) {$\la$};
                    \node[above=3pt] at (b) {$\lb$};
                    \node[below=3pt] at (c) {$\lc$};
                \end{scope}
            }
        % single edge
        \node[point] (x) at (8,0.7) {};
        \node[point] (y) at (10,0.7) {};
        \draw (x) -- (y);
        \node[below=3pt] at (x) {$1$};
        \node[below=3pt] at (y) {$2$};
    \end{tikzpicture}
    \caption{Left: two disjoint triangles (not connected). Right: a single edge (bipartite).}
    \label{fig:counter-0910}
\end{figure}

\begin{lemma}[Both Conditions Are Necessary]\label{lem:necessary-0910}
    Neither condition of \cref{thm:convergence-0910} can be dropped.
    \begin{enumerate}
        \item (Not connected.) Let $G$ be two disjoint triangles, on the vertices $1, 2, 3$ and
              $4, 5, 6$. This graph is not bipartite. The walk converges from every starting vertex, but the
              limit depends on the start, and the stationary distribution is not unique.
        \item (Bipartite.) Let $G$ be a single edge on the vertices $1, 2$. This graph is connected. The
              walk started at vertex $1$ does not converge, although $\pi = \left(\sfrac{1}{2}, \sfrac{1}{2}\right)$
              is stationary.
    \end{enumerate}
\end{lemma}

\begin{proof}
    \emph{Part 1.} The graph contains a triangle, which is an odd cycle, so it is not bipartite. There is
    no edge between $\set{1,2,3}$ and $\set{4,5,6}$, so $P_{uv} = 0$ whenever $u$ and $v$ are in different
    triangles. Start the token at vertex $1$. Then $p_t^4 = p_t^5 = p_t^6 = 0$ for all $t$, by induction on
    $t$: it holds for $t = 0$, and if it holds for $t$ then for $v \in \set{4,5,6}$ every term of
    $p_{t+1}^v = \sum_u p_t^u P_{uv}$ is $0$, because $p_t^u = 0$ for $u \in \set{4,5,6}$ and $P_{uv} = 0$ for
    $u \in \set{1,2,3}$. So the token stays in the first triangle. Every vertex of that triangle has degree $2$ in
    $G$, so the rows and columns $1, 2, 3$ of $P$ form exactly the transition matrix of the triangle, and
    the walk is the random walk on a triangle started at vertex $1$. By \cref{lem:triangle-0908},
    \[p_t \longrightarrow \left(\sfrac{1}{3}, \sfrac{1}{3}, \sfrac{1}{3}, 0, 0, 0\right) =: \pi^{(1)}.\]
    By the same argument the walk started at vertex $4$ converges to
    $\pi^{(2)} = \left(0, 0, 0, \sfrac{1}{3}, \sfrac{1}{3}, \sfrac{1}{3}\right)$, and by symmetry the walk from any
    other vertex converges to the uniform distribution on its own triangle. Since
    $\pi^{(1)} \neq \pi^{(2)}$, there is no single limit that is reached from every start. What fails here is
    the independence of the start, not the convergence.

    Both limits are stationary: $\pi^{(1)}$ is $0$ on the second triangle and uniform on the first, and the
    uniform distribution is stationary for the triangle, so $\pi^{(1)} P = \pi^{(1)}$; similarly for
    $\pi^{(2)}$. So the stationary distribution is not unique.

    \emph{Part 2.} A single edge is connected, and it is bipartite with the two sides $\set{1}$ and
    $\set{2}$. Both degrees are $1$, so
    \[P = \begin{pmatrix} 0 & 1 \\ 1 & 0 \end{pmatrix}, \qquad (a, b) P = (b, a).\]
    Starting from $p_0 = (1, 0)$ we get $p_t = (1, 0)$ for even $t$ and $p_t = (0, 1)$ for odd $t$, by
    induction on $t$. The first entry is the sequence $1, 0, 1, 0, \dots$, which has no limit, so $p_t$ does
    not converge. On the other hand $\left(\sfrac{1}{2}, \sfrac{1}{2}\right) P = \left(\sfrac{1}{2}, \sfrac{1}{2}\right)$,
    so $\pi = \left(\sfrac{1}{2}, \sfrac{1}{2}\right)$ is stationary (it is $\sfrac{\deg(v)}{2m}$ with $m = 1$): a
    stationary distribution exists, but the walk started at a vertex does not approach it.
\end{proof}

The second example shows what goes wrong in every bipartite graph: if all edges go between two sides
$L$ and $R$, then a walk started at a vertex of $L$ is in $L$ at all even times and in $R$ at all odd
times, so the probability $\sum_{v \in L} p_t^v$ of the side $L$ alternates between $1$ and $0$.

In terms of the eigenvalues, the two examples are exactly the two ways in which the proof in
\cref{sec:eigen-0910} can fail. If $G$ is bipartite with sides $L$ and $R$, let $x_u = 1$ for $u \in L$
and $x_u = -1$ for $u \in R$. All neighbours of a vertex are on the other side, so by
\cref{eq:average-0910} we have $Px = -x$: the number $-1$ is an eigenvalue, and its powers $(-1)^t$ do
not tend to $0$. For the single edge, $x = (1, -1)^{T}$. If $G$ is not connected, let $x_u = 1$ for the
vertices $u$ of one connected component and $x_u = 0$ for all others. All neighbours of a vertex are in
its own component, so $Px = x$, although $x$ is not a multiple of $\mathbf{1}$: the eigenvalue $1$ occurs
more than once, and the limit depends on the start. For the two triangles, $x = (1, 1, 1, 0, 0, 0)^{T}$.
